% Complete editorial arguments for the bounded corrected-author-source reconstruction.
\section{Domain convergence, approximation and recursive models}

\begin{annexitem}{annex:p10.domain.basis}{thm:p10-domain-basis}{Basis consequences and finite interpolation}
\begin{proof}
Testing \(x\ll y\) against the singleton directed set \(\{y\}\)
gives \(x\sqsubseteq y\). If \(x'\sqsubseteq x\ll y\sqsubseteq y'\)
and \(y'\sqsubseteq\bigsqcup^{\uparrow}A\), the witness for \(x\ll y\)
also bounds \(x'\); this proves the second implication.
Choose the directed subset \(C\subseteq B_x\) supplied by the basis
definition. Every member of \(B_x\) lies below \(x\), whereas
\(\bigsqcup^{\uparrow}C=x\). For \(b,b'\in B_x\), approximation
against \(C\) gives \(c,c'\in C\) above them, and directedness of
\(C\) gives a common upper bound in \(B_x\). Thus \(B_x\) is
directed with supremum \(x\). If \(x\) is compact, its approximation
against \(B_x\) gives \(b\in B_x\) with \(x\sqsubseteq b\sqsubseteq x\),
so \(x=b\in B\). Taking \(B=D\) proves the intrinsic continuous
criterion in both directions. If a basis consists of compact elements,
it contains all compact elements by the preceding argument, and hence
equals \(K(D)\); compactness below \(x\) is the same as compact
approximation to \(x\). This proves the algebraic criterion.

For the basis form of interpolation fix \(y\), and set
\[
 C_y=\{b\in B:(\exists c\in B)\ b\ll c\ll y\}.
\]
It is nonempty: choose \(c\in B_y\) and then \(b\in B_c\).
For \(b\ll c\ll y\) and \(b'\ll c'\ll y\), take
\(d\in B_y\) above \(c,c'\). Then \(b,b'\ll d\).
Approximation against the directed \(B_d\) and a common upper
bound in that set give \(a\in B_d\) above \(b,b'\).
Since \(a\ll d\ll y\), this \(a\) belongs to \(C_y\).
Thus \(C_y\) is directed. For every \(c\in B_y\),
\(B_c\subseteq C_y\), so its supremum is above \(c\).
All of \(C_y\) lies below \(y\); hence
\(\bigsqcup^{\uparrow}C_y=y\).
For each of the finitely many \(m\ll y\), take a member of
\(C_y\) above it and a common upper bound \(a\in C_y\).
Write \(a\ll b\ll y\), \(b\in B\). Then \(m\ll b\ll y\).
For empty \(M\), choose any \(b\in B_y\).
Finally, if \(c\ll\bigsqcup^{\uparrow}A\), interpolate
\(c\ll d\ll\bigsqcup^{\uparrow}A\), choose \(a\in A\) above
\(d\), and get \(c\ll a\). The reverse inclusion follows from
\(a\sqsubseteq\bigsqcup^{\uparrow}A\).
\end{proof}
\end{annexitem}

\begin{annexitem}{annex:p10.domain.examples}{thm:p10-domain-basic-examples}{Set information, partial graphs, interval approximation and a failed domain}
\begin{proof}
A directed union of subsets is their least upper bound. If a finite
\(M\subseteq B\) is covered by a directed family of subsets, choose
one family member containing each of its finitely many elements, then
one common upper member. Thus \(M\ll B\); for \(M=\varnothing\)
any member works. Conversely the finite subsets of \(B\) are directed
with union \(B\). If \(A\ll B\), one contains \(A\), so \(A\)
is finite and contained in \(B\). This proves algebraicity and the
exact compact-element description. For uncountable \(S\), any chain
of finite subsets has at most one distinct member of each finite
cardinality, because comparable finite sets of equal size are equal.
Its union is countable. It cannot be cofinal in all finite subsets
of \(S\), since cofinality would cover every singleton.

For a directed family of partial-function graphs, two assigned values
at the same input occur together in a common extending graph, so are
equal. Its union is therefore a partial function and is the least
upper bound. A finite graph is compact by the same finite-common-upper-
bound argument. Every graph is the directed union of its finite
subgraphs; if it is compact it must itself equal a finite subgraph.
The empty graph is bottom, including when either underlying set is
empty. This proves the partial-function assertion.

The real interval has the supremum of every nonempty subset. If
\(x<y\sqsubseteq\bigsqcup A\), some member of \(A\) is above
\(x\), or \(x\) would be an upper bound contradicting \(x<y\).
Zero is below every member of every nonempty directed family.
If \(x=y>0\), an increasing sequence strictly below \(x\) with
supremum \(x\) defeats approximation; if \(x>y\), the singleton
\(\{y\}\) defeats it. This establishes the exact relation and
compact elements. Rational approximants below each positive \(y\),
together with zero, form a directed set with supremum \(y\).
It is not algebraic because its only compact element cannot
approximate a positive point to its full value.

For the final example take elements \(\bot,a_n,b_n\ (n\in\mathbb N),\top\),
with two otherwise incomparable increasing chains, bottom below all,
and top above all. A directed set meeting both legs must contain top;
a directed set confined to a leg either has a greatest finite member
or is unbounded with supremum top. Thus it is a dcpo. The opposite
leg has supremum top and no member above a given leg element, so no
nonbottom element approximates any element above it. Top also fails
to approximate itself. Only bottom approximates every point, and
its singleton cannot have a nonbottom supremum. The dcpo is not
continuous.
\end{proof}
\end{annexitem}

\begin{annexitem}{annex:p10.domain.ideals}{thm:p10-domain-ideal-completion}{Round ideals, exact approximation and inverse recovery maps}
\begin{proof}
The empty finite-set interpolation clause makes \(i(b)\) nonempty;
transitivity makes it lower, and finite interpolation makes it directed
with respect to \(\prec\). A directed union of ideals is nonempty and
lower. Any two of its members lie in a common constituent ideal,
which supplies their strict upper member. Hence directed unions are
ideals and are least upper bounds.
For an ideal \(A\), roundness and lower closure give
\(A=\bigcup_{a\in A}i(a)\); this family is directed, since a
common strict upper member in \(A\) bounds two such ideals.
If \(b\in A'\) and a directed union of ideals contains \(A'\),
one constituent contains \(b\), hence contains \(i(b)\).
Thus \(i(b)\ll A'\), and so every \(A\subseteq i(b)\) is
way below \(A'\). Conversely, approximate \(A\ll A'\) against
the displayed directed family \(\{i(b):b\in A'\}\); this gives
the required \(b\). These observations prove both the relation and
that \(i(B)\) is a basis.
For reflexive \(\prec\), \(b\in i(b)\), so \(i(b)\) is
compact. Any compact ideal \(A\) is contained in some \(i(b)\),
\(b\in A\), and lower closure gives the opposite inclusion.
For a poset, \(b\mapsto\downarrow b\) is injective and reflects
order; hence these are exactly the compact elements.

Now let \(B\) be a basis of continuous \(D\), with relation \(\ll\).
Finite interpolation shows \(B_x\) is a round ideal, not merely
ordinary-order directed. If \(A\) is such an ideal, its supremum
\(x\) exists in \(D\). Each \(a\in A\) approximates a later
member of \(A\), and hence \(x\), so \(A\subseteq B_x\).
If \(b\in B_x\), interpolate \(b\ll c\ll x\) with \(c\in B\).
Approximation against \(A\) gives \(a\in A\) above \(c\),
so \(b\ll a\) and lower closure puts \(b\) in \(A\).
Thus \(B_x=A\). Conversely \(\bigsqcup^{\uparrow}B_x=x\).
These maps are inverse order isomorphisms, so preserve directed
suprema and are continuous. The final formula is just preservation
of the directed supremum \(x=\bigsqcup^{\uparrow}B_x\).
\end{proof}
\end{annexitem}

\begin{annexitem}{annex:p10.domain.topology}{thm:p10-domain-scott-topology}{Scott opens and the exact continuity equivalence}
\begin{proof}
Empty and whole are Scott-open. A union is upper, and if a directed
supremum belongs to it, it belongs to one open member and therefore
some directed-set member belongs to the union. For a binary intersection,
find a member of the directed family in each open and a common upper
member; upper closure puts the latter in both. Induction handles finite
intersections. Thus these sets form a topology.
Upper closure proves that \(x\sqsubseteq y\) implies the stated
neighborhood inclusion. If \(x\not\sqsubseteq y\), the complement
of \(\downarrow y\) is an open containing \(x\) and omitting \(y\):
\(\downarrow y\) is lower and contains the supremum of every
directed subset of itself. This proves the converse and \(T_0\).

For an order-continuous \(f\), inverse images of upper sets are
upper. If \(f(\bigsqcup^{\uparrow}A)\) belongs to an open, its
continuity and that open's inaccessibility put \(f(a)\) there for
some \(a\in A\). Thus inverse images are open.
Conversely suppose \(f\) is topologically continuous. Neighborhood
order recovery implies that it is monotone. Given directed \(A\),
put \(z=\bigsqcup^{\uparrow}f(A)\). The inverse image of the
closed \(\downarrow z\) is closed and contains \(A\), so contains
\(\bigsqcup^{\uparrow}A\). Hence
\(f(\bigsqcup^{\uparrow}A)\sqsubseteq z\); monotonicity gives
the reverse inequality. This proves the equivalence.
In a continuous domain, \(\mathord{\uparrow\!\uparrow}b\) is
upper. If \(b\ll\bigsqcup^{\uparrow}A\), interpolation gives
\(b\ll c\ll\bigsqcup^{\uparrow}A\), and some \(a\in A\)
above \(c\) belongs to this set. It is therefore open. For
\(x\in O\), approximate \(x\) by \(B_x\). Openness gives
\(b\in O\cap B_x\), whence
\(x\in\mathord{\uparrow\!\uparrow}b\subseteq O\).
This proves the topological-basis assertion.
\end{proof}
\end{annexitem}

\begin{annexitem}{annex:p10.domain.functions}{thm:p10-domain-function-space}{Pointwise convergence, joint evaluation and coordinate approximation}
\begin{proof}
For directed \(H\subseteq[D\longrightarrow E]\), define
\(h(x)=\bigsqcup^{\uparrow}_{f\in H}f(x)\). It is monotone.
For directed \(A\subseteq D\), the jointly indexed values
\(f(a)\), \((f,a)\in H\times A\), are directed. The same
upper bounds characterize the iterated suprema, so
\[
 h(\bigsqcup^{\uparrow}A)
 =\bigsqcup_{f\in H}^{\uparrow}\bigsqcup_{a\in A}^{\uparrow}f(a)
 =\bigsqcup_{a\in A}^{\uparrow}\bigsqcup_{f\in H}^{\uparrow}f(a)
 =\bigsqcup_{a\in A}^{\uparrow}h(a).
\]
Pointwise upper-bound comparison proves that \(h\) is the supremum
in the function order. Constants are continuous, so the bottom
assertion follows. For a directed family of strict maps the value
at bottom remains bottom, proving closure of the strict space.
Coordinatewise suprema give products, continuous projections and
continuous pairing. If \(x\ll x'\), \(y\ll y'\), test a directed
set of pairs whose supremum bounds \((x',y')\); witnesses for each
coordinate have a common upper pair, proving approximation.
Conversely test with directed pairs \((a,y')\) for the first
coordinate and \((x',b)\) for the second. Product bases follow by
forming products of the directed coordinate approximants.

For a directed set \(C\subseteq[D\longrightarrow E]\times D\),
with coordinate families \(H,A\),
\[
 (\bigsqcup^{\uparrow}H)(\bigsqcup^{\uparrow}A)
          =\bigsqcup_{f\in H,\,a\in A}^{\uparrow} f(a)
          =\bigsqcup_{(f,a)\in C}^{\uparrow} f(a).
\]
For the last equality a common upper member of \(C\) bounds any
two members supplying the two chosen coordinates. Evaluation is
monotone, and this proves its continuity on the actual product.
For continuous \(k:D\times E\to F\), each \(k(d,-)\) is
continuous. The map \(d\mapsto k(d,-)\) preserves directed
suprema pointwise by continuity of \(k\) on the fixed-second-
coordinate family. Its inverse uses continuous evaluation and
pairing. Both composites are identity by evaluation at every pair.
The operations taking a function to its curry or uncurry preserve
directed suprema pointwise. Composition is continuous because at
each argument its directed-supremum calculation is the same joint-
evaluation calculation, with a common upper pair replacing two
indices. This also proves the functorial action and its local
continuity, with reversed composition in its first argument.
\end{proof}
\end{annexitem}

\begin{annexitem}{annex:p10.domain.fixpoint}{thm:p10-domain-least-fixpoint}{The least value, its continuous operator and a recursive graph calculation}
\begin{proof}
Since \(\bot\sqsubseteq f(\bot)\), monotonicity makes
\(f^n(\bot)\) an increasing chain. Put its supremum \(d\).
Continuity gives \(f(d)=\bigsqcup_n^{\uparrow}f^{n+1}(\bot)=d\):
deleting the first, least member does not change the upper bounds.
If \(f(x)\sqsubseteq x\), induction gives
\(f^n(\bot)\sqsubseteq x\), and hence \(d\sqsubseteq x\).

Let \(I_n(f)=f^n(\bot)\). These maps are monotone in \(f\):
inductively, \(f\sqsubseteq g\) and \(I_n(f)\sqsubseteq I_n(g)\)
imply \(f(I_n(f))\sqsubseteq g(I_n(g))\).
They are continuous by induction. For directed \(H\), assuming
the assertion at \(n\),
\[
 I_{n+1}(\bigsqcup^{\uparrow}H)
   =\bigsqcup_{g,f\in H}^{\uparrow}g(I_n(f))
   =\bigsqcup_{f\in H}^{\uparrow}f(I_n(f)).
\]
The first equality uses pointwise suprema and each \(g\)'s
continuity; the second uses a common upper \(h\in H\) for
\(f,g\), giving \(g(I_n(f))\sqsubseteq h(I_n(h))\).
The \(I_n\)'s increase pointwise; their supremum in the function
dcpo is therefore continuous and is exactly \(\operatorname{fix}\).
Currying the continuous \(F:P\times D\to D\) and composing
with this operator proves the parameter assertion and its equation.

For the factorial instance, graph extension of \(p\) extends
\(\Phi(p)\). At a directed union, any defined output at positive
\(n\) uses one input graph pair and is therefore already produced
at some member; the zero output is present at every member.
Thus \(\Phi\) preserves directed graph unions. Starting with the
empty graph, induction shows the \(k\)-th graph consists exactly
of \((n,n!)\) for \(n<k\). The union is total factorial, and
the theorem makes it the least fixed point. Identity on the partial-
function dcpo has every point fixed and its least fixed point is
bottom, establishing the contrasting divergent equation.
\end{proof}
\end{annexitem}

\begin{annexitem}{annex:p10.domain.induction}{thm:p10-domain-fixpoint-principles}{Admissible induction and strict uniformity}
\begin{proof}
By induction \(f^n(\bot)\in P\), starting at its required bottom
member and using \(f(P)\subseteq P\). Closure under this increasing
chain's supremum puts \(\operatorname{fix}(f)\) in \(P\).
For uniformity, strictness starts the equality
\(h(f^0(\bot_D))=g^0(\bot_E)\). The commuting equation gives
\(h(f^{n+1}(\bot_D))=g(h(f^n(\bot_D)))=g^{n+1}(\bot_E)\).
Continuity of \(h\) carries the supremum of this chain to the
supremum of the chain on the right. This is exactly the claimed
fixed-point equality, and exhibits the separate role of strictness.
\end{proof}
\end{annexitem}

\begin{annexitem}{annex:p10.domain.failures}{thm:p10-domain-failed-premises}{Explicit missing-premise counterexamples}
\begin{proof}
Every directed subset of a discrete poset is a singleton, so every
function on it is continuous. Successor on \(\mathbb N\) has no
fixed point and its domain has no bottom.
The ordinal chain \(\omega+2\) is a pointed dcpo. The stated
function is monotone in all finite, limit and top cases. Its finite
iterates from zero are the finite naturals, with supremum \(\omega\),
but its value there is \(\omega+1\). It does have fixed point
\(\omega+1\); the counterexample defeats the omega-iteration
formula without continuity, not the separate monotone existence
statement. It explicitly fails continuity on the directed finite
naturals.
On \(\omega+1\), extended successor preserves directed suprema:
a directed set with finite supremum has its greatest member, and
one with supremum \(\omega\) either contains that point or has
unbounded finite members whose successors also have supremum
\(\omega\). Its least fixed point is \(\omega\). The finite
predicate contains zero and is successor-closed but not closed
under this chain's supremum, so cannot be used for admissible
induction.
Finally let \(D=\{*\}\), \(E=\{0<1\}\), let \(f,g\) be
their identities, and set \(h(*)=1\). This continuous non-strict
map satisfies \(hf=gh\), but
\(h(\operatorname{fix}(f))=1\ne0=\operatorname{fix}(g)\).
\end{proof}
\end{annexitem}

\begin{annexitem}{annex:p10.domain.constructors}{thm:p10-domain-constructors}{Directed sets in lifted, tagged and smashed domains}
\begin{proof}
A directed subset of a lifting is either its new-bottom singleton,
or has a nonempty directed old part which is cofinal in it.
Its supremum is respectively new bottom or the lift of the old
supremum. New bottom is compact, and
\(\operatorname{up}(b)\ll\operatorname{up}(x)\) holds exactly
when \(b\ll x\) in \(D\), as follows by testing old directed
sets and by discarding the non-cofinal new-bottom member.
This proves its stated bases and continuity/algebraicity preservation.
Extending \(f:D\to E\) by new bottom \(\mapsto\bot_E\)
is continuous by this classification and is strict. Restriction
to old elements is the inverse operation. The bijection is natural,
since restriction and extension commute with composition of the
correctly typed maps.

In the coalesced sum, a directed set containing a nonbottom member
has all its nonbottom members in the same branch: two different
branches have no common upper bound. The nonbottom part is cofinal,
and its supremum is computed in that branch. Otherwise the
supremum is bottom. Approximation of nonbottom tagged elements
is exactly their branch approximation; bottom is compact.
Tagged bases, together with bottom, therefore give the required
bases. Two strict maps out of the factors agree at their common
bottom and uniquely paste to a strict map on the sum. That map
is continuous by the directed-set classification. Conversely
restriction gives the two strict maps. Without strictness they
can disagree at bottom, for example the two constant maps with
values zero and one into \(\{0<1\}\), so the asserted coproduct
property is not a property of all pointed continuous maps.

For smash product a directed set having a nonbottom member has a
cofinal family of nonbottom pairs; its two coordinates are directed,
and their suprema remain nonbottom. Its supremum is that pair.
The remaining case is the bottom singleton. Testing with directed
pairs varying one coordinate, and using a common upper pair for
the converse, proves coordinatewise approximation for nonbottom
pairs. Their paired bases plus bottom are bases; if a coordinate's
basis approximants initially include bottom, its nonbottom part
is cofinal whenever the approximated coordinate is nonbottom.
This proves both preservation assertions.
A strict map from smash product corresponds exactly to a
continuous map on cartesian product which is bottom on either
bottom axis. Currying identifies these with strict maps from
\(D\) into the strict function dcpo on \(E,F\). Currying and
uncurrying remain inverse and continuous by the function-space
proof. The monoidal unit is \(I_{\bot}\), the two-element chain:
its unique nonbottom element paired with \(d\ne\bot\) identifies
\(I_{\bot}\otimes D\) with \(D\). Associativity and symmetry
are the maps on nonbottom triples and swapped pairs, extended
strictly at bottom. They preserve the classified directed suprema;
their coherence equations are the ordinary tuple equations on
nonbottom elements and the bottom equation in the other cases.
\end{proof}
\end{annexitem}

\begin{annexitem}{annex:p10.domain.steps}{thm:p10-domain-algebraic-functions}{Finite observations form an actual compact function basis}
\begin{proof}
In a pointed bounded-complete dcpo every bounded subset has a
supremum: its finite joins, including the empty join bottom, form
a directed family with exactly the upper bounds of that subset.
For compact \(a\), the set \(\uparrow a\) is Scott-open by
the compact-element definition. Thus the stated step function
is continuous. It lies below a continuous \(f\) exactly when
\(b\sqsubseteq f(a)\). It is compact in the function dcpo:
if a directed family has supremum above it, compactness of
\(b\) at argument \(a\) gives a family member whose value at
\(a\) bounds \(b\); monotonicity makes that member bound
the whole step.

Consider finitely many steps with a pointwise upper bound. Their
pointwise join \(s\) is monotone. For a directed \(A\subseteq D\),
each step active at \(\bigsqcup^{\uparrow}A\) has its compact
input threshold below some member of \(A\). One common upper
member activates all those finitely many thresholds. No inactive
step can become active below the supremum. Hence at that member
\(s\) already has its value at the supremum, proving continuity.
If a directed function family has supremum above \(s\), obtain
one witness above each step and a common upper member; it bounds
their join. Thus \(s\) is compact.
For any \(f\), all finite joins of steps below it are directed,
including the empty join, and are pointwise bounded by \(f\).
For compact \(b\sqsubseteq f(x)\), continuity gives
\(f(x)=\bigsqcup_{a\in K(D)\cap\downarrow x}^{\uparrow}f(a)\).
Compactness supplies such an \(a\) with \(b\sqsubseteq f(a)\).
Its step is below \(f\) and at \(x\) is above \(b\).
Since \(E\) is algebraic, these observations give \(f\) itself
as the supremum of the finite-step joins. Any compact \(f\)
must equal one of these joins, so the basis characterization is
exact, not only a list of some compact functions.

For a bounded family \(H\) of functions let \(u\) be a continuous
upper bound and set \(h(x)=\bigsqcup_{f\in H}f(x)\). The required
suprema exist below \(u(x)\); an empty family gives bottom.
For directed \(A\), all values \(f(a)\) are bounded by
\(u(\bigsqcup^{\uparrow}A)\). Equality of upper bounds permits
interchanging their two suprema:
\[
 h(\bigsqcup^{\uparrow}A)
   =\bigsqcup_{f\in H}\bigsqcup_{a\in A}^{\uparrow}f(a)
   =\bigsqcup_{a\in A}^{\uparrow}\bigsqcup_{f\in H}f(a).
\]
Thus \(h\) is continuous and is the least upper function, proving
bounded completeness. Finite subsets of a countable set of compact
input/output pairs are countable, so the displayed compact basis
is countable under the stated countability premise.
\end{proof}
\end{annexitem}

\begin{annexitem}{annex:p10.domain.bilimit}{thm:p10-domain-bilimit}{Commuting tuples, both mediators and the limit basis}
\begin{proof}
For an embedding--projection pair, \(e(x)\sqsubseteq z\) iff
\(x\sqsubseteq p(z)\): apply \(p\) in one direction and use
\(ep\sqsubseteq1\) in the other. Thus \(e\) is a lower
adjoint and \(p\) its unique upper adjoint. Between pointed
objects \(e(\bot)\sqsubseteq ep(\bot)\sqsubseteq\bot\), so
\(e\) is strict; applying \(p\) to \(\bot\sqsubseteq e(\bot)\)
shows \(p\) is strict too. Composed adjunctions give exactly
the claimed equations for the \(p_{nm}\)'s.

Coordinatewise directed suprema of commuting tuples commute,
because each \(p_{nm}\) is continuous. Their bottom tuple
commutes by strictness, so \(D\) is a pointed dcpo.
The displayed \(e_nx\) is a commuting tuple: on coordinates
below \(n\) use composition of projections, above \(n\) use
\(p_{ij}e_{jn}=e_{in}\), and across \(n\) use
\(p_{ij}e_{jn}=p_{in}\). These identities follow from
\(p_{ij}e_{ji}=1\) and the connecting-map laws.
Each coordinate is continuous, so \(e_n\) is continuous;
\(p_n\) is continuous by coordinatewise convergence.
Directly \(p_ne_n=1\). For a commuting tuple \(x\), its
\(i\le n\) coordinates are unchanged by \(e_np_n\), and
its \(i\ge n\) coordinates become
\(e_{in}p_{ni}(x_i)\sqsubseteq x_i\). Thus
\(e_np_n\sqsubseteq1\). Moreover
\(e_n=e_me_{mn}\), \(p_n=p_{nm}p_m\), so
\(e_np_n\sqsubseteq e_mp_m\) for \(n\le m\).
Every fixed coordinate equals its original value once \(n\)
reaches that coordinate, proving \(\bigsqcup_n^{\uparrow}e_np_n=1\).

A projection cone \((g_n:C\to D_n)\) gives its unique map
\(g(c)=(g_n(c))_n\) into the tuple dcpo; it is continuous
coordinatewise. For an embedding cocone \((h_n:D_n\to E)\),
the maps \(h_np_n\) increase, since
\(h_np_n=h_me_{mn}p_{nm}p_m\sqsubseteq h_mp_m\).
Their pointwise supremum \(h\) is continuous. For \(n\ge m\),
\(p_ne_m=e_{nm}\), so \(he_m=h_m\).
Indeed a cofinal tail of the defining supremum is constantly
\(h_ne_{nm}=h_m\). If \(k\) is another mediator,
\[
 k=k\bigsqcup_n^{\uparrow}e_np_n
       =\bigsqcup_n^{\uparrow}ke_np_n
       =\bigsqcup_n^{\uparrow}h_np_n=h.
\]
This proves both universal properties using actual maps.

For a cocone of pairs \((u_n,v_n)\), uniqueness of upper
adjoints to its embeddings gives \(v_n=p_{nm}v_m\).
Consequently \(v:E\to D\), \(v(x)=(v_nx)_n\), is a cone
mediator, while \(u:D\to E\) is the colimit mediator
\(\bigsqcup_n^{\uparrow}u_np_n\).
For \(m\ge n\), \(v_nu_m=p_{nm}v_mu_m=p_{nm}\).
Hence the \(n\)-th coordinate of \(vu(x)\) is
\(\bigsqcup_{m\ge n}^{\uparrow}p_{nm}p_m(x)=x_n\), proving
\(vu=1_D\). Also
\(uv=\bigsqcup_n^{\uparrow}u_np_nv
=\bigsqcup_n^{\uparrow}u_nv_n\).
Thus the supremum condition makes these inverse maps and makes
the cocone colimiting. Conversely a colimiting cocone has the
unique isomorphism to the constructed colimit, with its unique
projection mates, so transports the identity supremum to itself.

An embedding with continuous projection preserves and reflects
\(\ll\). For preservation, if \(e(y)\sqsubseteq\bigsqcup A\),
apply \(p\), use its continuity, and use \(x\ll y\) to
get \(x\sqsubseteq p(a)\), hence \(e(x)\sqsubseteq a\).
For reflection test \(e(x)\ll e(y)\) on \(e(A)\), and
cancel the order embedding.
For \(x\in D\) let
\(C_x=\bigcup_ne_n((B_n)_{p_nx})\).
Its members approximate \(x\). Two of them can be transported
to a common later stage \(k\); each still approximates \(p_kx\),
so directedness of \((B_k)_{p_kx}\) supplies a common upper
approximant there. Thus \(C_x\) is directed. Its supremum is
\[
 \bigsqcup_n^{\uparrow}e_n
       (\bigsqcup^{\uparrow}(B_n)_{p_nx})
         =\bigsqcup_n^{\uparrow}e_np_nx=x.
\]
This proves the basis assertion. If the stages are algebraic,
the displayed basis consists of compact elements by preservation
of \(\ll\). Any compact \(x\) must equal one member of
the directed compact approximation family \(C_x\), proving
the exact formula for \(K(D)\). The final countability statement
uses a countable union of the stated countable bases.
\end{proof}
\end{annexitem}

\begin{annexitem}{annex:p10.domain.mixed}{thm:p10-domain-mixed-embedding}{Adjoint directions, functor laws and preservation of the identity supremum}
\begin{proof}
For input pairs \((e,p)\), \((e',p')\), functoriality gives
\[
 F(e,p')F(p,e')=F(pe,p'e')=1,
 \qquad
 F(p,e')F(e,p')=F(ep,e'p')\sqsubseteq F(1,1)=1.
\]
The inequality uses monotonicity on correctly typed hom objects,
which is part of Scott-continuity. Thus the maps are an e-p-pair.
If another pair follows, the projection of the composite is the
reversed composite of projections. Contravariance reverses it
back in \(F\)'s first argument, giving the covariant functor
law for \(\widehat F\); identities give its identity law.

Take synchronized expanding sequences with bilimits \(D,D'\)
and pairs \((e_n,p_n),(e'_n,p'_n)\). The images
\(k_n=F(p_n,e'_n)\), \(q_n=F(e_n,p'_n)\) form a cocone
of embedding--projection pairs. Their approximating idempotents are
\[
 k_nq_n=F(e_np_n,e'_np'_n).
\]
The two argument families are increasing with supremum identity.
Local continuity gives the supremum over all pairs of indices;
the diagonal is cofinal because every two indices have a common
upper index. Hence \(\bigsqcup_n^{\uparrow}k_nq_n=F(1,1)=1\).
The already proved bilimit criterion makes this cocone colimiting.
For identical synchronized sequences this also proves the diagonal
functor's continuity.
For function spaces the action is \(F(f,g)(h)=ghf\).
Composition and evaluation in the function-space proof give
\[
 (\bigsqcup_i^{\uparrow}g_i)h(\bigsqcup_j^{\uparrow}f_j)
       =\bigsqcup_{i,j}^{\uparrow}g_i h f_j
\]
pointwise, because \(h\) and each \(g_i\) are continuous.
This verifies the required hypothesis rather than assuming
function-space bilimit preservation from its name.
\end{proof}
\end{annexitem}

\begin{annexitem}{annex:p10.domain.fold}{thm:p10-domain-recursive-solution}{Tail cocones and the explicit inverse fold/unfold maps}
\begin{proof}
The maps \(G^n(a)\) are embeddings because \(G\) is a functor
on the embedding category. Thus the already proved bilimit exists.
Discarding the first stage does not change its colimit: a cocone
over the tail uniquely extends to stage zero by composing its
first leg with \(a\), and the same assertion holds for mediators.
Therefore \((D,(e_{n+1})_n)\) is a colimit of \((D_{n+1})_n\).
Continuity of \(G\) says exactly that
\((G(D),(G(e_n))_n)\) is another colimit of that same diagram.
Their two unique mediators are inverse: either composite preserves
all colimit legs and thus equals the identity by uniqueness.

The explicit colimit-mediator formula of the bilimit theorem gives
\(\operatorname{unfold}=\bigsqcup_n^{\uparrow}G(e_n)p_{n+1}\).
Using the projection mates of the cocone \(G(e_n)\) gives
\(\operatorname{fold}=\bigsqcup_n^{\uparrow}e_{n+1}G(e_n)^*\).
These are increasing families by the cocone/pair equations, not
suprema of arbitrary unproved compatible expressions. The leg
equation is \(\operatorname{unfold}e_{n+1}=G(e_n)\);
taking its unique upper adjoints also gives
\(G(e_n)^*=p_{n+1}\operatorname{fold}\).
For the one-point seed the bottom selection and unique map back
are an e-p-pair, so this particular sequence is available without
an additional seed-existence assumption. No uniqueness among
all possible equation solutions was used or proved by this argument.
\end{proof}
\end{annexitem}

\begin{annexitem}{annex:p10.domain.recursive-examples}{thm:p10-domain-recursive-examples}{Prefix bilimits and a nontrivial reflexive object}
\begin{proof}
In prefix order, two words with a common upper word are comparable.
Thus any directed family is a compatible family of prefixes. Its
union is a word, finite or infinite, and is its least upper bound.
A finite prefix of that union is contained in one constituent:
take its finitely many positions and a common upper constituent.
So finite words are compact. An infinite word is the supremum
of its finite prefixes, with none containing it, and is not compact.
Each word is the directed supremum of its finite prefixes, proving
algebraicity. The empty word is bottom.
Removing the first letter of a nonempty word and inserting it
again give inverse maps for the displayed lifted-product equation.
Both preserve order. As inverse order isomorphisms of dcpo's,
they preserve directed suprema, so are continuous.

Let \(W_n\) be words of length at most \(n\). Truncation
\(t_n\) and inclusion \(j_n\) satisfy
\(t_nj_n=1\), \(j_nt_n\sqsubseteq1\).
Truncation preserves a directed union: if the union has at least
\(n\) letters, one constituent already contains that prefix;
otherwise the finite union itself occurs in a constituent.
Thus these are continuous e-p-pairs.
Their truncation idempotents increase with supremum identity,
so the bilimit criterion identifies the word dcpo with their
unfolding limit. The order isomorphisms
\(W_{n+1}\cong(A\times W_n)_{\bot}\), with \(W_0=I\),
identify the actual adjacent embeddings with successive liftings
of the bottom seed embedding. This also proves the asserted
canonical-unfolding identification, not just the existence of
some unrelated word model.

The flat \(A=\mathbb N_{\bot}\) is an algebraic pointed
bounded-complete dcpo: a directed set either has only bottom or
has one nonbottom number as its largest member, and a bounded
pair has its larger comparable member. Constants give a continuous
map \(a:A\to[A\longrightarrow A]\) pointwise; evaluation at
bottom is continuous by the function theorem. We have
\(pa=1_A\), and \((ap)(f)\) is the constant \(f(\bot)\)
function, below \(f\) by monotonicity. This is an actual e-p-pair.
The preceding mixed-variance theorem applies to function space
and supplies its continuous embedding functor; the fold theorem
therefore supplies the inverse maps of the claimed equation.
The seed embeds into the bilimit, and has distinct natural-number
elements, so the resulting object cannot be one-point. At every
stage the finite-step theorem gives a pointed bounded-complete
algebraic function domain; hence the bilimit is algebraic by
the bilimit basis theorem. All those compact bases are countable,
so it is also omega-algebraic.
For seed \(I\), \([I\longrightarrow I]=I\) and every connecting
map is identity, so its tuple bilimit is exactly one-point.

Finally define application \(D\times D\to D\) as evaluation
after \(\operatorname{unfold}\) in the first coordinate.
For a continuous contextual body \(b:\Gamma\times D\to D\),
abstraction is \(\operatorname{fold}\circ\Lambda b\).
Applying it to an argument evaluates \(\Lambda b\), since
\(\operatorname{unfold}\operatorname{fold}=1\); this gives beta.
Conversely currying application to a contextual \(d:\Gamma\to D\)
gives \(\operatorname{unfold}d\) by transpose uniqueness, and
folding gives \(d\), proving eta. Pairing, evaluation and currying
make these actual continuous maps. This proves the stated semantic
equations without asserting an unproved operational adequacy or
full-abstraction theorem.
\end{proof}
\end{annexitem}
